% SOURCE_AUTHORITY: sga5_fr_workpass.tex, lines 6549--6600 (LF-normalized unit).
% SOURCE_SHA256: 791F4EFFC5E02832D5D77ED03518C8156D6F07E4C8238B03545DB93D883FBB28
\begin{statement}{Lema 6.23.3}
(cf. SGA $4^{1/2}$ Rapport 4.7). Sean $P$ un $G_V$-módulo localmente factor directo de un módulo libre de tipo finito, $P'=i_!P$, y $Q$ un complejo de $\Lambda[G]$-módulos, perfecto sobre $\Lambda$. Entonces $P\lten_{\Lambda}Q$, considerado como objeto de $D({}_{\Lambda[G]}V)$ mediante la acción diagonal de $G$, es perfecto sobre $\Lambda[G]$, y se tiene un cuadrado conmutativo
\begin{equation}
\tag{*}
\begin{tikzcd}[column sep=huge,row sep=large]
\delta^*\uRHom_{\Lambda[G]}(p_1^*P',p_2^!P')\lten_{\Lambda}
\uRHom_{\Lambda[G]}(Q,Q) \arrow[r,"(1)"] \arrow[d,"(2)"']
& K_Y \arrow[d,equal]\\
\delta^*\uRHom_{\Lambda[G]}(p_1^*P'\lten_{\Lambda}Q,p_2^!P'\lten_{\Lambda}Q) \arrow[r,"(3)"']
& K_Y ,
\end{tikzcd}
\end{equation}
donde $\delta:Y\to Y\times_SY$ es la diagonal, la flecha horizontal superior es el producto tensorial de las flechas
\[
\delta^*\uRHom_{\Lambda[G]}(p_1^*P',p_2^!P')\xrightarrow{6.5.3}
K\Lambda[G]_{Y,\natural}\xrightarrow{x\mapsto x(e)}K_Y
\]
y
\[
\uRHom_{\Lambda[G]}(Q,Q)\xrightarrow{\Tr_{\Lambda}}\Lambda,
\]
la flecha vertical se deduce del producto tensorial externo por restricción de escalares mediante la diagonal
\[
\Lambda[G]\to\Lambda[G]\otimes\Lambda[G],
\]
y la flecha horizontal inferior es la compuesta de 6.5.3 (definida porque $P'\lten_{\Lambda}Q\in\ob D_{\mathrm{ctf}}({}_{\Lambda[G]}Y)$) y de $x\mapsto x(e)$\footnote{Debería poder sustituirse $P$ por un complejo perfecto (sobre $\Lambda[G]$) (e incluso $P'$ por un objeto de $D_{\mathrm{ctf}}({}_{\Lambda[G]}Y)$).}.
\end{statement}

La primera afirmación es evidente, pues se reduce al caso $P=\Lambda[G]_V$, y entonces $P\lten_{\Lambda}Q$ es isomorfo al producto tensorial de $P$ por $Q$ considerado con la acción trivial de $G$. Para la segunda, obsérvese primero que, según 6.5.4 y 6.5.5, se tiene
\[
\delta^*\uRHom_{\Lambda[G]}(p_1^*P',p_2^!P')
\xrightarrow{\sim}
D_{\Lambda[G]}(i_!P)\lten_{\Lambda[G]} i_!P
\xrightarrow{\sim}
i_!(D_{\Lambda[G]}P\lten_{\Lambda[G]}P),
\]
de modo que se reduce a comprobar que el cuadrado inducido por (*) sobre $V$ es conmutativo. Por tanto, se puede suponer que $V=Y$. El cuadrado (*) se escribe entonces
\begin{equation}
\tag{**}
\begin{tikzcd}[column sep=huge,row sep=large]
\uRHom_{\Lambda[G]}(P,K_Y\lten P)\lten\uRHom_{\Lambda[G]}(Q,Q)
  \arrow[r,"(1)"] \arrow[d,"(2)"']
& K_Y \arrow[d,equal]\\
\uRHom_{\Lambda[G]}(P\lten Q,K_Y\lten P\lten Q) \arrow[r,"(3)"']
& K_Y ,
\end{tikzcd}
\end{equation}
donde (2) es el producto tensorial externo, (1) viene dada por
\[
\Tr_{\Lambda[G]}(-)(e)\otimes\Tr_{\Lambda},
\]
(trazas en el sentido del n\textsuperscript{o}~5), y (3) por $\Tr_{\Lambda[G]}(-)(e)$. Por descenso, se reduce al caso $P=\Lambda[G]_Y$, y, por la observación inicial, se puede suponer entonces que $G$ actúa trivialmente sobre $Q$. La afirmación es, en tal caso, consecuencia inmediata de 5.12.5.
